% © 2026 Ing. David Jaroš — CC BY-NC-ND 4.0
%
% This work is licensed under a Creative Commons Attribution-NonCommercial-NoDerivatives
% 4.0 International License (CC BY-NC-ND 4.0).
% See LICENSE.md for full license text.

% UBT-AI-PROVENANCE-BEGIN
% schema: ubt-ai-provenance/v1
% tier: C_working
% ai_assistance: disclosed
% human_review: risk-based
% editorial_responsibility: Ing. David Jaroš
% policy: ../../AI_PROVENANCE.md
% notice: Working material; exhaustive human review is not claimed.
% UBT-AI-PROVENANCE-END

\documentclass[11pt,a4paper]{article}
\usepackage{amsmath,amssymb,amsthm}
\usepackage{booktabs}
\usepackage{geometry}
\usepackage{hyperref}
\geometry{margin=2.5cm}

\newtheorem{proposition}{Proposition}
\newtheorem{remark}{Remark}
\newtheorem{obstruction}{Obstruction}

\newcommand{\proved}[1]{\textbf{[PROVED]} #1}
\newcommand{\cond}[1]{\textbf{[CONDITIONAL]} #1}
\newcommand{\heur}[1]{\textbf{[HEURISTIC]} #1}
\newcommand{\nogo}[1]{\textbf{[NO-GO]} #1}

\title{T-Dual Winding Correction Audit for the $B$-Coefficient}
\author{Ing.\ David Jaro\v{s}\\
\small \texttt{research\_tracks/rg\_B46/}}
\date{May 2026}

\begin{document}
\maketitle

\begin{abstract}
We test whether the winding correction often written numerically as
$\Delta B_{\mathrm{wind}}\approx 18.5$ is derivable from the UBT action, or must be
formally downgraded. We enforce hard constraints: no fitted inputs, no fixed
special integer input, and no target-driven tuning. The result is: (i) the
self-dual normalization $R_{\psi}=1$ is not derived unconditionally from the
current UBT moduli sector, (ii) T-duality of the compact $\psi$ spectrum is
exact at one-loop Gaussian level when momentum and winding sectors are both
included, (iii) the winding correction is $n$-dependent in all explicit
non-fitted formulas, and (iv) it does not supply a universal constant shift of
the coefficient of $n\log n$.
\end{abstract}

\tableofcontents

\section{Setup and target statement}

The baseline potential form used in this track is
\begin{equation}
V_{\mathrm{eff}}(n)=n^2-B\,n\log n.
\end{equation}
The target claim under audit is a radius/self-duality based winding correction
\begin{equation}
\Delta B_{\mathrm{wind}}\stackrel{?}{=}\text{constant}.
\end{equation}

We test five items:
\begin{enumerate}
\item derive $R_{\psi}=1$ from UBT moduli,
\item prove T-duality symmetry of compact $\psi$ sector,
\item compute winding contribution without fixed special integer input,
\item decide constant vs $n$-dependent structure,
\item decide whether winding modifies $B$ (coefficient of $n\log n$) or adds a separate term.
\end{enumerate}

\section{Radius selection from the UBT moduli sector}

From the compact $(t_E,\psi)$ quadratic sector used in the prior self-dual
torus derivation, one has
shape and scale variables
\begin{equation}
\rho=\frac{R_t}{R_{\psi}},\qquad s=\sqrt{R_tR_{\psi}},
\end{equation}
and one-loop free energy stationarity at $\rho=1$ under isotropic assumptions.

\begin{proposition}[What is actually derived]
In the current compact-cycle analysis, square-shape stationarity
$R_t=R_{\psi}$ is conditional, while the overall scale $s$ is not fixed by the
kinetic-plus-mass sector alone.
\end{proposition}

\begin{proof}
The cited derivation gives $\partial_{\rho}F|_{\rho=1}=0$ and local shape
stability, but explicitly leaves $\partial_sF=0$ without an isolated finite
solution unless extra physics is added.
\end{proof}

\begin{obstruction}
Without an independent scale-fixing equation, one cannot promote
$R_t=R_{\psi}$ to an unconditional prediction $R_{\psi}=1$.
\end{obstruction}

\cond{Radius-normalization statement: one may set $R_{\psi}=1$ by unit choice, but
this is normalization, not a derived modulus-fixing theorem.}

\section{T-duality test in compact $\psi$ sector}

For a compact bosonic direction with momentum index $k\in\mathbb Z$ and winding
index $w\in\mathbb Z$, the Gaussian one-loop lattice sum is
\begin{equation}
Z(R_{\psi};t)=\sum_{k,w\in\mathbb Z}
\exp\!\left[-\pi t\left(\frac{k^2}{R_{\psi}^2}+w^2R_{\psi}^2\right)\right].
\end{equation}
Under
\begin{equation}
R_{\psi}\mapsto \frac{1}{R_{\psi}},\qquad k\leftrightarrow w,
\end{equation}
$Z$ is invariant (Poisson resummation / lattice relabelling).

\begin{proposition}[Spectral T-duality]
At one-loop Gaussian level, with both momentum and winding sectors retained, the
compact-spectrum partition sum is invariant under $R_{\psi}\leftrightarrow 1/R_{\psi}$.
\end{proposition}

\begin{remark}
This is a spectral statement. Extending it to full UBT action covariance needs
modular-weight and measure-compensation inputs that remain open in the modular
covariance audit.
\end{remark}

\cond{T-duality status for this task: valid in the compact Gaussian sector;
not yet a full-action theorem.}

\section{Winding contribution without fixed special integer input}

Using the one-loop KK threshold structure already used in this track,
\begin{equation}
B_{\mathrm{mom}}(n)=\frac{N_{\mathrm{eff}}}{12\pi^2}\,n,
\end{equation}
with $N_{\mathrm{eff}}=12$ from field-content audit. If T-dual symmetry is imposed
at the level of equal momentum/winding weighting, then
\begin{equation}
B_{\mathrm{mom+wind}}(n)=2B_{\mathrm{mom}}(n)
=\frac{N_{\mathrm{eff}}}{6\pi^2}\,n,
\end{equation}
and therefore
\begin{equation}
\Delta B_{\mathrm{wind}}(n)=B_{\mathrm{mom+wind}}(n)-B_{\mathrm{mom}}(n)
=\frac{N_{\mathrm{eff}}}{12\pi^2}\,n.
\label{eq:deltaBwind_general}
\end{equation}

Equation \eqref{eq:deltaBwind_general} uses no fixed special integer input and
contains no fit parameter.

\section{Constant vs $n$-dependent classification}

\begin{proposition}
The winding correction derived from the compact momentum/winding tower is
$n$-dependent, not a universal constant.
\end{proposition}

\begin{proof}
Equation \eqref{eq:deltaBwind_general} is linear in $n$ for fixed
$N_{\mathrm{eff}}$. Hence $\Delta B_{\mathrm{wind}}$ varies across KK levels and
cannot define one universal number independent of $n$.
\end{proof}

\nogo{Any standalone constant claim $\Delta B_{\mathrm{wind}}\equiv\text{const}$ is not
obtained from the explicit non-fitted tower calculation.}

\section{Does winding modify $B$ of $n\log n$ or add a separate term?}

If one inserts $B(n)=B_0+\Delta B_{\mathrm{wind}}(n)$ into
$V_{\mathrm{eff}}=n^2-B(n)n\log n$, one gets
\begin{equation}
V_{\mathrm{eff}}(n)
=n^2-B_0n\log n
-\frac{N_{\mathrm{eff}}}{12\pi^2}n^2\log n.
\end{equation}
Thus the winding piece has $n^2\log n$ scaling and is not a universal shift of a
constant coefficient multiplying $n\log n$.

\begin{proposition}
In the explicit tower derivation, winding enters as an additional
$n$-dependent structure. It does not provide a universal constant renormalization
of the $n\log n$ coefficient.
\end{proposition}

\section{Verdict}

\begin{center}
\begin{tabular}{p{0.56\linewidth}p{0.32\linewidth}}
\toprule
Requirement & Outcome \\
\midrule
Derive $R_{\psi}=1$ from UBT moduli & \textbf{CONDITIONAL} (shape yes, scale not fixed) \\
Prove T-duality of compact $\psi$ sector & \textbf{CONDITIONAL} (spectral Gaussian level) \\
Compute winding term without fixed special integer & \textbf{PROVED}: Eq.\ \eqref{eq:deltaBwind_general} \\
Constant vs $n$-dependent & \textbf{PROVED}: $n$-dependent \\
Shift of constant $B$ for $n\log n$ & \textbf{NO-GO}: induces separate $n^2\log n$ structure \\
\bottomrule
\end{tabular}
\end{center}

\bigskip
\noindent\textbf{Final class for the claim ``constant winding correction is derived''}:\
\begin{equation}
\boxed{\textbf{NO-GO}}.
\end{equation}

\noindent The previous constant-number usage is therefore downgraded to
\textbf{HEURISTIC bookkeeping at a chosen evaluation level}, not a derived UBT
constant.

\end{document}
